% A template showing Latex possibilities for providing datasheets
% Initially provided by smartfpga for the fpga-systems.ru community

\documentclass[a4paper, 10pt, onecolumn]{article}


\newcommand {\documentname}{LaTeX documentation demonstrator template}
\newcommand {\documenthdrstring}{\textit{\textbf{\documentname}}}


\usepackage[left=2cm,right=2cm,top=2cm,bottom=2cm,bindingoffset=0cm]{geometry}

\usepackage[utf8]{inputenc}        

\usepackage[english]{babel} 
\usepackage{url}                	% optional for Internet Links

\usepackage{helvet}
\renewcommand{\familydefault}{\sfdefault}
\fontfamily{phv}\selectfont

\usepackage{boldline}
\usepackage{tabularx}
\usepackage[
singlelinecheck=false % <-- important
]{caption}

\usepackage{colortbl}               % color in tables
\usepackage{color}


\definecolor{dimared}{RGB}{106,42,50}
\definecolor{gray50}{gray}{0.5}
\definecolor{gray25}{gray}{0.75}
\definecolor{lightgray}{gray}{0.9}

\usepackage{graphicx}
\usepackage{lscape}

\usepackage{amsmath}                % optional for formulas
\usepackage{float}
%\usepackage[font=sf]{floatrow}

\usepackage{sectsty}
\allsectionsfont{\normalfont\sffamily\bfseries} 

%\usepackage{bytefield}


\usepackage{fancyhdr}
\pagestyle{fancy}
\fancypagestyle{title}{%
    \fancyhead{}
    \renewcommand{\headrulewidth}{0pt}
    \rhead{\includegraphics[width=4cm]{fig/fpgasys.png}}
    %\lhead{\ipcorename}
}
\setlength\headheight{46pt} %% just to make warning go away. Adjust the value after looking into the warning.
\rhead{\includegraphics[width=4cm]{fig/fpgasys.png}}
\lhead{\documenthdrstring}
\textheight = 695pt

\def\widehline{%
\noalign{\global\dimen1 \arrayrulewidth
\global\arrayrulewidth3\arrayrulewidth}%
\hline
\noalign{\global\arrayrulewidth\dimen1 }}

\newcommand{\chline}{%
    \arrayrulecolor{dimared}\widehline\arrayrulecolor{gray25}
}


\usepackage[
    colorlinks=true,
    urlcolor=blue,
    linkcolor=black,
    citecolor=black
]{hyperref}

\begin{document}
\thispagestyle{title}
\arrayrulecolor{gray25}

\begin{flushleft}
\huge \documentname
\normalsize
\end{flushleft}

\section*{Revision History}
The following table shows the revision history for this document. But also it shows how a full-width table may look like.

\begin{table}[H]
\label{tab:revhist}
\centering\bgroup
\def\arraystretch{1.5}
\begin{tabularx}{\textwidth}{|l|l|l|X|}
    \hline
    \textbf{Date} & \textbf{Version} & \textbf{Author} & \textbf{Revision} \\ \chline
    20.04.2021 & 1.0 & Dmitry Eliseev & The LaTeX template for your datasheets.\\ \hline
\end{tabularx}
\egroup
\end{table}


\section*{General Description}
This document shows possibilities of LaTeX for creating datasheets. 
It is available as a Github-project at \url{https://github.com/smartfpga/Latex-Datasheet-Template.git} \cite{GithubLink}.


\section*{Insert figures}
Please put your figures in the /fig directory. Here comes Figure \ref{fig:FigureRef}.
\begin{figure}[H]
    \begin{center}
        \includegraphics[width=14cm]{fig/SampleFigure.png}
    \end{center}
    \caption{Caption of an example figure.}
    \label{fig:FigureRef}
\end{figure}



% no bibtex is used for the bibliography. You should take care yourself about formatting the bibliography entries. 
\begin{thebibliography}{3}
\bibitem{GithubLink}
Gitlab Project: LaTeX Documentation Template. URL:\\
\url{https://github.com/smartfpga/Latex-Datasheet-Template.git}
\end{thebibliography}

\end{document}
